\documentclass[11pt]{article}
\newcommand{\DocNumber}{VOLUME 05}
\newcommand{\DocTitle}{Secure C++23 Engineering Guide}
\newcommand{\DocSubtitle}{Object safety, concurrency and input trust}
\newcommand{\DocSummary}{Twelve security-control modules, defensible CWE distinctions, a review scorecard and adversarial labs.}
% Shared, pdfLaTeX-only design system for the C++23 documentation suite.
\usepackage[T1]{fontenc}
\usepackage[utf8]{inputenc}
\usepackage{lmodern}
\usepackage[english]{babel}
\usepackage[letterpaper,margin=0.82in,headheight=16pt,headsep=17pt]{geometry}
\usepackage{microtype}
\usepackage{amsmath,amssymb}
\usepackage{array,booktabs,tabularx,longtable}
\usepackage[table]{xcolor}
\usepackage[most]{tcolorbox}
\usepackage{listings}
\usepackage{fancyhdr}
\usepackage{titlesec}
\usepackage{enumitem}
\usepackage{hyperref}
\usepackage{needspace}
\usepackage{etoolbox}
\definecolor{ink}{HTML}{13253B}
\definecolor{teal}{HTML}{006E7F}
\definecolor{muted}{HTML}{5A6C7C}
\definecolor{fog}{HTML}{EDF4F6}
\definecolor{sand}{HTML}{FAF3E8}
\definecolor{rule}{HTML}{C9D7DF}
\definecolor{codebg}{HTML}{F7F9FB}
\definecolor{codecomment}{HTML}{47715D}
\hypersetup{colorlinks=true,linkcolor=teal,urlcolor=teal,pdfauthor={Independent C++23 engineering documentation},pdftitle={C++23 Technical Suite}}
\setlength{\parindent}{0pt}
\setlength{\parskip}{6pt}
\setcounter{tocdepth}{2}
\renewcommand{\arraystretch}{1.18}
\titleformat{\section}{\Large\bfseries\color{ink}}{\thesection}{0.7em}{}
\titleformat{\subsection}{\large\bfseries\color{teal}}{\thesubsection}{0.55em}{}
\titlespacing*{\section}{0pt}{17pt}{7pt}
\titlespacing*{\subsection}{0pt}{11pt}{4pt}
\pagestyle{fancy}
\fancyhf{}
\fancyhead[L]{\footnotesize\bfseries\color{teal}C++23 / ISO 14882:2024}
\fancyhead[R]{\footnotesize\color{muted}\DocNumber}
\fancyfoot[L]{\scriptsize\color{muted}Independent engineering companion, not an ISO publication}
\fancyfoot[R]{\scriptsize\color{muted}\thepage}
\renewcommand{\headrulewidth}{0.35pt}
\lstdefinestyle{cppstyle}{language=C++,basicstyle=\ttfamily\footnotesize,keywordstyle=\color{teal}\bfseries,commentstyle=\color{codecomment},stringstyle=\color{ink},backgroundcolor=\color{codebg},frame=single,rulecolor=\color{rule},breaklines=true,breakatwhitespace=true,columns=fullflexible,keepspaces=true,showstringspaces=false,tabsize=2,numbers=none,upquote=true,xleftmargin=3pt,xrightmargin=3pt,aboveskip=6pt,belowskip=8pt}
\lstset{style=cppstyle}
\newtcolorbox{standardbox}{enhanced,breakable,colback=fog,colframe=teal!35!white,boxrule=0.55pt,arc=1.4mm,left=3mm,right=3mm,top=2mm,bottom=2mm,before skip=7pt,after skip=8pt}
\newtcolorbox{practicebox}{enhanced,breakable,colback=sand,colframe=ink!20!white,boxrule=0.55pt,arc=1.4mm,left=3mm,right=3mm,top=2mm,bottom=2mm,before skip=7pt,after skip=8pt}
\newcommand{\anchorref}[1]{\par\smallskip{\footnotesize\textbf{Standard locator:} \texttt{#1}}\par}
\newcommand{\standard}[1]{\begin{standardbox}\textbf{Standard-grounded interpretation.} #1\end{standardbox}}
\newcommand{\practice}[1]{\begin{practicebox}\textbf{Engineering recommendation (not a normative ISO rule).} #1\end{practicebox}}
\newcommand{\takeaway}[1]{\par\textbf{Review question.} #1\par}
\newcommand{\docfront}{%
\begin{titlepage}
\vspace*{25mm}
{\color{teal}\bfseries\large ENGINEERING REFERENCE SERIES\par}
\vspace{12mm}
{\raggedright\color{ink}\bfseries\fontsize{24}{30}\selectfont\DocTitle\par}
\vspace{7mm}
{\color{muted}\fontsize{15}{19}\selectfont\DocSubtitle\par}
\vspace{12mm}
\noindent\color{teal}\rule{\linewidth}{2pt}\par
\vspace{9mm}
{\large\DocSummary\par}
\vfill
\begin{tabularx}{\linewidth}{@{}lX@{}}
\textbf{Volume} & \DocNumber\\
\textbf{Standards basis} & ISO/IEC 14882:2024 (seventh edition, C++23)\\
\textbf{Document type} & Original interpretive and instructional material\\
\textbf{Date} & 16 October 2026\\
\textbf{Toolchain} & pdfLaTeX; illustrative C++23 examples require separate compiler support\\
\end{tabularx}
\vspace{8mm}
{\footnotesize This document is not an authorized copy of, replacement for, or amendment to ISO/IEC 14882:2024. The original standard controls in the event of disagreement.}\par
\end{titlepage}
\pagenumbering{roman}
\tableofcontents
\clearpage
\pagenumbering{arabic}
\section*{How to use this volume}\addcontentsline{toc}{section}{How to use this volume}
The reference is \emph{ISO/IEC 14882:2024, Programming languages --- C++}. Section locators and square-bracketed stable names identify where to verify precise wording; they are not quotations. \textbf{Standard-grounded interpretation} boxes summarize language or library semantics. \textbf{Engineering recommendation} boxes contain optional design, security, portability, or process guidance. Code is original and illustrative; individual library facilities may not yet be implemented by every vendor. This suite does not grant rights to reproduce or distribute the licensed standard.\par\medskip
}
\newcommand{\docend}{\clearpage\section*{Reference and traceability}\addcontentsline{toc}{section}{Reference and traceability}
\textbf{Primary source.} ISO/IEC 14882:2024(en), \emph{Programming languages --- C++}, seventh edition, October 2024. Use its numbered clauses, stable names, Annex C (informative compatibility), and Annex D (compatibility features) for authoritative details.\par
\textbf{Limits.} Clause-level summaries and exercises in this document are selective explanations, not exhaustive reproductions of all semantic constraints, library mandates, complexity clauses, and overload participation requirements. A successful compiler test demonstrates behavior of that specific implementation and configuration, not universal conformance.\par
\textbf{Use of contemporary engineering tools.} Compilers, sanitizers, analyzers, CI systems, repository policies, and vulnerability taxonomies are implementation guidance rather than language-standard requirements.\par}

\begin{document}
\docfront
\section{Assurance boundary}
The C++ language standard defines what programs mean or when it imposes no requirements; it does not certify that a program is secure. This document introduces optional security controls and selected Common Weakness Enumeration (CWE) mappings as modern engineering additions. CWE mappings are examples, not an exhaustive or normative crosswalk.

\small
\begin{longtable}{@{}p{.18\linewidth}p{.40\linewidth}p{.34\linewidth}@{}}
\toprule
\textbf{Source construct} & \textbf{Failure condition} & \textbf{Control objective}\\
\midrule
\endfirsthead
\toprule
\textbf{Source construct} & \textbf{Failure condition} & \textbf{Control objective}\\
\midrule
\endhead
Raw pointer/reference & Referent dead or inaccessible & Explicit owner + lifetime proof\\
\addlinespace[2pt]
Index arithmetic & Value out of valid range & Checked bounds and overflow\\
\addlinespace[2pt]
Shared mutable state & Missing happens-before edge & Synchronization discipline\\
\addlinespace[2pt]
Non-owning view & Owner changed or destroyed & Borrowing contract\\
\addlinespace[2pt]
Exception path & Partially updated invariant & Commit/rollback structure\\
\addlinespace[2pt]
File access & Untrusted path or race & OS-backed authorization\\
\addlinespace[2pt]
\bottomrule
\end{longtable}
\normalsize
\section{C++ semantic classification and threat modeling}
\anchorref{3; 4.1; 6.9}
A secure coding rule can be stricter than the language standard. Undefined behavior, ill-formed code, implementation-defined choices, and data races must be handled according to their actual semantic categories.

\standard{The referenced C++ clauses define the language and library semantics that can underlie this weakness; they do not prescribe a cybersecurity control program.}
\begin{itemize}[leftmargin=*,itemsep=1pt,topsep=3pt]
\item Identify adversarial inputs and preconditions.
\item Create a minimally unsafe test as a negative fixture, isolated from production.
\item Test the repaired implementation with appropriate instrumentation.
\end{itemize}
\begin{lstlisting}[language=C++]
// Review checklist, not a complete program:
// Untrusted input -> parse -> validated domain type -> operation
// Record ownership, bounds, integer domain, and failure mode.
\end{lstlisting}
\practice{Create a threat model for trust boundaries and each resource lifetime; do not equate clean compilation with safe execution.}
\takeaway{Security evidence target: Evidence: input trust map, threat assumptions and violated precondition catalog.}
\section{Out-of-bounds access and container misuse}
\anchorref{6.7; 7.6; 24.2--24.3}
Indexing raw storage or a vector outside its valid range is not prevented by the type system. at() can provide an explicit checked access contract.

\standard{The referenced C++ clauses define the language and library semantics that can underlie this weakness; they do not prescribe a cybersecurity control program.}
\begin{itemize}[leftmargin=*,itemsep=1pt,topsep=3pt]
\item Identify adversarial inputs and preconditions.
\item Create a minimally unsafe test as a negative fixture, isolated from production.
\item Test the repaired implementation with appropriate instrumentation.
\end{itemize}
\begin{lstlisting}[language=C++]
#include <vector>
#include <optional>
std::optional<int> lookup(const std::vector<int>& v,
                          std::size_t index) {
    if (index >= v.size()) return std::nullopt;
    return v[index];
}
\end{lstlisting}
\practice{Keep range checks at the trust boundary and avoid unchecked narrowing from signed input.}
\takeaway{Security evidence target: CWE-125 / CWE-787 are useful external taxonomy examples, not ISO C++ rule identifiers.}
\section{Use-after-free and dangling references}
\anchorref{6.7.3--6.7.7; 20; 23.3}
Object lifetime and view lifetime are independent. A valid-looking string\_view can still reference destroyed or reallocated data.

\standard{The referenced C++ clauses define the language and library semantics that can underlie this weakness; they do not prescribe a cybersecurity control program.}
\begin{itemize}[leftmargin=*,itemsep=1pt,topsep=3pt]
\item Identify adversarial inputs and preconditions.
\item Create a minimally unsafe test as a negative fixture, isolated from production.
\item Test the repaired implementation with appropriate instrumentation.
\end{itemize}
\begin{lstlisting}[language=C++]
#include <string>
#include <string_view>
std::string safe_name() {
    std::string owned{"item"};
    return owned; // return owning object
}
// Never return string_view referring to a local string.
\end{lstlisting}
\practice{Do not persist borrowed views unless the owner and invalidation conditions are explicit.}
\takeaway{Security evidence target: CWE-416 and CWE-825 provide external classification vocabulary.}
\section{Integer overflow, conversion and allocation math}
\anchorref{6.8; 7.3--7.4; 7.6}
Signed arithmetic overflow can be undefined; conversions, narrowing and size arithmetic create hazards even when behavior is otherwise defined.

\standard{The referenced C++ clauses define the language and library semantics that can underlie this weakness; they do not prescribe a cybersecurity control program.}
\begin{itemize}[leftmargin=*,itemsep=1pt,topsep=3pt]
\item Identify adversarial inputs and preconditions.
\item Create a minimally unsafe test as a negative fixture, isolated from production.
\item Test the repaired implementation with appropriate instrumentation.
\end{itemize}
\begin{lstlisting}[language=C++]
#include <limits>
#include <optional>
#include <cstddef>
std::optional<std::size_t> add_sizes(std::size_t a,
                                      std::size_t b) {
    if (b > std::numeric_limits<std::size_t>::max() - a)
        return std::nullopt;
    return a + b;
}
\end{lstlisting}
\practice{Perform precondition arithmetic in a domain that cannot overflow; use bounded allocations.}
\takeaway{Security evidence target: CWE-190 / CWE-681 are external labels; include adversarial tests near numeric limits.}
\section{Uninitialized or indeterminate data}
\anchorref{6.7; 9.4}
Using values before permitted initialization can have undefined behavior in specified contexts; the treatment of particular uninitialized-friendly types and operations requires close reading of the applicable clauses.

\standard{The referenced C++ clauses define the language and library semantics that can underlie this weakness; they do not prescribe a cybersecurity control program.}
\begin{itemize}[leftmargin=*,itemsep=1pt,topsep=3pt]
\item Identify adversarial inputs and preconditions.
\item Create a minimally unsafe test as a negative fixture, isolated from production.
\item Test the repaired implementation with appropriate instrumentation.
\end{itemize}
\begin{lstlisting}[language=C++]
struct Settings {
    bool enabled{false};
    unsigned retries{0};
};
Settings s{};
\end{lstlisting}
\practice{Use definite initialization, default member initializers and value-initialization for scalar state.}
\takeaway{Security evidence target: Require initialization before each control-flow read.}
\section{Type punning, alignment and object representation}
\anchorref{6.7; 7.6; 20}
Object representation, alignment, aliasing, and lifetime rules constrain access to bytes reinterpreted as other types.

\standard{The referenced C++ clauses define the language and library semantics that can underlie this weakness; they do not prescribe a cybersecurity control program.}
\begin{itemize}[leftmargin=*,itemsep=1pt,topsep=3pt]
\item Identify adversarial inputs and preconditions.
\item Create a minimally unsafe test as a negative fixture, isolated from production.
\item Test the repaired implementation with appropriate instrumentation.
\end{itemize}
\begin{lstlisting}[language=C++]
#include <bit>
#include <cstdint>
#include <limits>
static_assert(sizeof(float) == sizeof(std::uint32_t));
std::uint32_t bits(float f) { return std::bit_cast<std::uint32_t>(f); }
\end{lstlisting}
\practice{Prefer bit\_cast for same-size trivially copyable bit representations and serialization for external formats.}
\takeaway{Security evidence target: Check the target representation and intended bit semantics, not merely that sizes happen to match.}
\section{Iterator and view invalidation}
\anchorref{24--26}
Container mutation can invalidate references, iterators or pointers. Lazy views can defer dereference until after the backing container changes.

\standard{The referenced C++ clauses define the language and library semantics that can underlie this weakness; they do not prescribe a cybersecurity control program.}
\begin{itemize}[leftmargin=*,itemsep=1pt,topsep=3pt]
\item Identify adversarial inputs and preconditions.
\item Create a minimally unsafe test as a negative fixture, isolated from production.
\item Test the repaired implementation with appropriate instrumentation.
\end{itemize}
\begin{lstlisting}[language=C++]
#include <vector>
void append_then_access(std::vector<int>& v, int x) {
    v.push_back(x);
    if (!v.empty()) { int y = v.back(); (void)y; }
}
\end{lstlisting}
\practice{Document invalidation behavior per container and mutation; avoid deferred closures over mutating storage.}
\takeaway{Security evidence target: Test capacity growth, erase, rehash and move transitions.}
\section{Concurrency, races and memory ordering}
\anchorref{6.9.2; 33}
Conflicting unsynchronized non-atomic accesses can create a data race and undefined behavior. Atomic operations need the right ordering for the intended communication protocol.

\standard{The referenced C++ clauses define the language and library semantics that can underlie this weakness; they do not prescribe a cybersecurity control program.}
\begin{itemize}[leftmargin=*,itemsep=1pt,topsep=3pt]
\item Identify adversarial inputs and preconditions.
\item Create a minimally unsafe test as a negative fixture, isolated from production.
\item Test the repaired implementation with appropriate instrumentation.
\end{itemize}
\begin{lstlisting}[language=C++]
#include <mutex>
struct State {
    std::mutex mu;
    int version{};
};
void change(State& s) {
    std::lock_guard lock{s.mu};
    ++s.version;
}
\end{lstlisting}
\practice{Prefer mutexes for compound shared invariants, and explicitly document happens-before edges.}
\takeaway{Security evidence target: CWE-362 is an external vulnerability category; run ThreadSanitizer independently of AddressSanitizer.}
\section{Exception safety and transactional updates}
\anchorref{14; 20; 24}
An exception can interrupt a function between state changes. The basic or strong exception guarantee depends on invariants and commit structure, not syntactic exception handling.

\standard{The referenced C++ clauses define the language and library semantics that can underlie this weakness; they do not prescribe a cybersecurity control program.}
\begin{itemize}[leftmargin=*,itemsep=1pt,topsep=3pt]
\item Identify adversarial inputs and preconditions.
\item Create a minimally unsafe test as a negative fixture, isolated from production.
\item Test the repaired implementation with appropriate instrumentation.
\end{itemize}
\begin{lstlisting}[language=C++]
#include <vector>
void replace_if_valid(std::vector<int>& target,
                      const std::vector<int>& proposed) {
    std::vector<int> candidate = proposed; // may throw
    target.swap(candidate);                  // commit
}
\end{lstlisting}
\practice{Build a prospective state first, then commit it with a nonthrowing operation when possible.}
\takeaway{Security evidence target: Inject allocation failures or throwing test elements and assert invariants.}
\section{Untrusted parsing and expensive computation}
\anchorref{23; 27; 31--32}
Parsing errors, unbounded allocation, and pathological regular expressions can exhaust resources without requiring a memory violation.

\standard{The referenced C++ clauses define the language and library semantics that can underlie this weakness; they do not prescribe a cybersecurity control program.}
\begin{itemize}[leftmargin=*,itemsep=1pt,topsep=3pt]
\item Identify adversarial inputs and preconditions.
\item Create a minimally unsafe test as a negative fixture, isolated from production.
\item Test the repaired implementation with appropriate instrumentation.
\end{itemize}
\begin{lstlisting}[language=C++]
#include <string_view>
#include <expected>
enum class InputError { too_long, empty };
std::expected<std::string_view, InputError>
validate(std::string_view x) {
    if (x.empty()) return std::unexpected(InputError::empty);
    if (x.size() > 4096)
        return std::unexpected(InputError::too_long);
    return x;
}
\end{lstlisting}
\practice{Apply input-size limits, parser state budgets and timeouts outside the language contract.}
\takeaway{Security evidence target: CWE-400 is an external label for resource exhaustion.}
\section{Filesystem paths and external process boundaries}
\anchorref{31.12; 31.10}
filesystem path normalization is not by itself a sandbox. Symbolic links, time-of-check/time-of-use and OS permissions affect the actual resource resolved.

\standard{The referenced C++ clauses define the language and library semantics that can underlie this weakness; they do not prescribe a cybersecurity control program.}
\begin{itemize}[leftmargin=*,itemsep=1pt,topsep=3pt]
\item Identify adversarial inputs and preconditions.
\item Create a minimally unsafe test as a negative fixture, isolated from production.
\item Test the repaired implementation with appropriate instrumentation.
\end{itemize}
\begin{lstlisting}[language=C++]
#include <filesystem>
#include <system_error>
namespace fs = std::filesystem;
std::error_code ec;
auto exists = fs::exists(fs::path{"data.txt"}, ec);
// An existence check is not a safe authorization decision.
\end{lstlisting}
\practice{Enforce permissions and trusted-root policies with appropriate OS primitives; treat paths as untrusted data.}
\takeaway{Security evidence target: CWE-22 and CWE-367 are external taxonomies, not guarantees from std::filesystem.}
\section{Dependency and artifact integrity}
\anchorref{16; 20--33}
Library correctness does not prove the security of your third-party implementation or its supply chain. Toolchains, standard libraries and transitive dependencies also carry risks.

\standard{The referenced C++ clauses define the language and library semantics that can underlie this weakness; they do not prescribe a cybersecurity control program.}
\begin{itemize}[leftmargin=*,itemsep=1pt,topsep=3pt]
\item Identify adversarial inputs and preconditions.
\item Create a minimally unsafe test as a negative fixture, isolated from production.
\item Test the repaired implementation with appropriate instrumentation.
\end{itemize}
\begin{lstlisting}[language=C++]
// Review evidence:
// compiler + libstdc++/libc++ identity
// dependency lockfile + artifact digest
// SAST/DAST/fuzz coverage + disclosed residual risk
\end{lstlisting}
\practice{Use pinned dependencies, SBOM/provenance evidence, vulnerability review and reproducible builds as appropriate.}
\takeaway{Security evidence target: Record remediation owner and service-level objectives outside ISO conformity.}
\section{Security review scorecard}
\small
\begin{longtable}{@{}p{.18\linewidth}p{.40\linewidth}p{.34\linewidth}@{}}
\toprule
\textbf{Control} & \textbf{Evidence required} & \textbf{Disposition}\\
\midrule
\endfirsthead
\toprule
\textbf{Control} & \textbf{Evidence required} & \textbf{Disposition}\\
\midrule
\endhead
Ownership & Owner identified for each externally visible view & Pass / action / exception\\
\addlinespace[2pt]
Bounds & All untrusted indices checked & Pass / action / exception\\
\addlinespace[2pt]
Integers & Range proof for lengths, allocation and offsets & Pass / action / exception\\
\addlinespace[2pt]
Concurrency & Thread-sharing model, lock protocol and races & Pass / action / exception\\
\addlinespace[2pt]
Failures & Exception guarantees and error propagation & Pass / action / exception\\
\addlinespace[2pt]
Parsers & Size, nesting, CPU and memory limits & Pass / action / exception\\
\addlinespace[2pt]
Build & Reproducible flags, sanitizer logs and dependencies & Pass / action / exception\\
\addlinespace[2pt]
\bottomrule
\end{longtable}
\normalsize
\section{Adversarial lab cases}
\begin{itemize}[leftmargin=*,itemsep=1pt,topsep=3pt]
\item Pass index -1 through a signed-to-size\_t conversion boundary and verify controlled rejection.
\item Invalidate vector storage through reserve/push\_back and confirm no borrowed references survive.
\item Fuzz a bounded record parser with zero-length, maximum-length and malformed encodings.
\item Run shared-state mutation under ThreadSanitizer and compare to a locked version.
\item Inject a throwing value into vector construction and verify the original committed state is preserved.
\end{itemize}
\docend
\end{document}
